% Plantilla para Tesis, Trabajo Fin de Máster y Trabajo Fin de Grado en el Departamento de Electrotecnia de la ETSII-UPM.
% Version 1.3.2 (2016-11-07)
% Por Hugo Rocha Mendonça
%
% Advertencia: Esta versión está en periodo de pruebas y puede contener errores.
%
% Para comunicar los fallos encontrados o para cualquier otra sugerencia relacionada con
% la plantilla, dirigirse a:
%
%             hugo.rocha@upm.es


\documentclass[tfm]{dinel-upm}
% \documentclass[<Documento>, <Impresion>, <Interlineado>]{dinel-upm}
% opciones del documento:
% <Documento>:
% 	tesis - doctorado
%	tfm - trabajo fin de máster
%	tfg - trabajo fin de grado (por defecto)
% <Impresion>:
%	unacara
%	doblecara (por defecto)
% <Interlineado>:
%	espaciosimple
%	espaciounoymedio (por defecto)
%	espaciodoble


%************************************************************
% PAQUETES EXTRAS


\usepackage{blindtext} % Genera texto automático


%************************************************************


%************************************************************
% DATOS DEL TRABAJO

\author{Nombre Completo} % nombre completo
\TitAutor{} % título del autor, si hay
\Director{Nombre Completo} % nombre completo del director del trabajo
\TitDirector{Título del director} % título del director
\title{Título del Trabajo} % título del trabajo
\codirector{Nombre Completo} % codirector
\TitCodirector{Título del director} % título del codirector
\degree{Máster Universitario en Automática y Robótica}
\authorname{Tu Nombre Completo}
%************************************************************


%************************************************************
% GENERA LISTAS DE SÍMBOLOS Y ACRÓNIMOS

\makeindex

\makenomenclature
% Ejecutar este código en el terminal en la misma carpeta del fichero principal:
% makeindex <filename>.nlo -s nomencl.ist -o <filename>.nls

\makeglossaries
% Ejecutar este código en el terminal en la misma carpeta del fichero principal:
% makeindex -s <filename>.ist -t <filename>.alg -o <filename>.acr <filename>.acn
%************************************************************


%************************************************************
% GENERA BIBLIOGRAFIA

\bibliography{biblio}
%************************************************************


%************************************************************
% EMPIEZA EL DOCUMENTO

\begin{document}

%------------------------------------------------------------
% Primera hoja:
\titulo
%------------------------------------------------------------
% Hoja del tribunal (sólo para la TESIS)
%\tribunal{Presidente}{Secretario}{Vocal 1}{Vocal 2}{Vocal 3}
%------------------------------------------------------------
% Dedicatoria y agradecimientos:
\dedicatoria{primeras_hojas/dedicatoria.tex}
\agradecimientos{primeras_hojas/agradecimientos.tex}
%------------------------------------------------------------
% Resúmenes:
% Para TESIS
\resumen{primeras_hojas/resumen.tex}
\abstract{primeras_hojas/abstract.tex}
% Para TFM y TFG
\resumenejecutivo{primeras_hojas/resumen_ejecutivo.tex}
%------------------------------------------------------------
% Indices general, de figuras y de tablas
\indices
%------------------------------------------------------------
% Lista de símbolos y de acrónimos (opcional):
% ¡OJO! Seguir las instrucciones arriba en: GENERA LISTAS DE SÍMBOLOS Y ACRÓNIMOS.
\simbolos
\acronimos
%------------------------------------------------------------

% Texto principal:
\texto

% Sugerencia: escribir los capítulos en ficheros aparte y incluirlos en el texto principal tal como enseñado a continuación:
\chapter{Introducción}
\Blindtext

\section{Sección 1 de la introducción}
\Blindtext

\section{Sección 2 de la introducción}
\Blindtext

\subsection{Subsección 2.1 de la introducción}
\blindtext

\subsubsection{Subsubsección 2.1.1 de la introducción}
\blindtext

\subsubsection{Subsubsección 2.1.2 de la introducción}
\blindtext

\subsection{Subsección 2.2 de la introducción}
\Blindtext

\section{Sección 3 de la introducción}
\Blindtext


\chapter{Objetivos}
\begin{enumerate}
    \item Objetivo 1
    \item Objetivo 2
    ...
\end{enumerate}
\include{capitulos/Metodologia}
\chapter{Resultados y discusión}
\Blindtext

\section{Valoración de impactos y de aspectos de responsabilidad legal, ética y profesional}
Si procede alrededor de 2000 palabras.
\chapter{Conclusión 26}
\Blindtext

\chapter{Líneas futuras - Opcional}
\Blindtext

% BIBLIOGRAFIA
\chapter{Bibliografía}
\printbibliography[heading=none]

\appendix
\chapter{Planificación temporal}
\Blindtext

\chapter{Presupuesto}
\Blindtext

\end{document}
%************************************************************ 